\chapter{Lesson 9: Studying Chinese}

\section{Vocabulary}

\begin{center}
\begin{tabular}{llll}
\toprule
Characters & Pinyin & Part of speech & Meaning \\
\midrule
\zh{学习} & xuéxí & verb & to study \\
\zh{复习} & fùxí & verb & to review \\
\zh{预习} & yùxí & verb & to preview \\
\zh{课} & kè & noun / MW & class, lesson \\
\zh{第} & dì & prefix & (ordinal) \\
\zh{语法} & yǔfǎ & noun & grammar \\
\zh{生词} & shēngcí & noun & new words \\
\zh{汉字} & Hànzì & noun & Chinese character \\
\zh{发音} & fāyīn & noun & pronunciation \\
\zh{写} & xiě & verb & to write \\
\zh{说} & shuō & verb & to speak, to say \\
\zh{教} & jiāo & verb & to teach \\
\zh{懂} & dǒng & verb & to understand \\
\zh{得} & de & particle & (links verb and complement) \\
\zh{慢} & màn & adjective & slow \\
\zh{快} & kuài & adjective & fast \\
\zh{容易} & róngyì & adjective & easy \\
\zh{难} & nán & adjective & difficult \\
\zh{太\ldots{}了} & tài \ldots le & & too \ldots; so \ldots \\
\zh{真} & zhēn & adverb & really \\
\zh{怎么} & zěnme & question word & how \\
\zh{早} & zǎo & adjective & early \\
\zh{晚} & wǎn & adjective & late \\
\zh{录音} & lùyīn & noun & recording \\
\zh{作业} & zuòyè & noun & homework \\
\zh{一点儿} & yìdiǎnr & & a little \\
\zh{有(一)点儿} & yǒu (yì)diǎnr & adverb & somewhat (negative) \\
\zh{平常} & píngcháng & adverb & usually \\
\bottomrule
\end{tabular}
\end{center}

\section{Ordinal numbers}

\zh{第} + number gives an ordinal: \zh{第一} ``first'', \zh{第二课} ``Lesson 2'',
\zh{第三个学生} ``the third student''. Note that \zh{一} keeps its first tone here:
\emph{dì-yī}.

\section{Complements of degree}

The single most important new grammar of the term. To \emph{evaluate} how an action is
done (well, fast, early \ldots), Chinese puts an adjective \emph{after} the verb, linked
by \zh{得} \emph{de}.

\begin{definition}[Complement of degree]
\begin{center}
Subject + Verb + \zh{得} + (adverb) + Adjective
\end{center}
\end{definition}

\begin{itemize}
  \item \zh{他说得很快。} He speaks fast.
  \item \zh{你写得真好！} You write really well!
  \item \zh{我昨天睡得很晚。} I went to bed late last night.
  \item \zh{她来得很早。} She came early.
\end{itemize}

\subsection{With an object: repeat the verb}

\zh{得} must directly follow the verb, so an object has to be moved out of the way.
Either repeat the verb, or put the object first (and drop the first verb).

\begin{center}
\begin{tikzpicture}[w/.style={draw, rounded corners, minimum height=0.8cm, align=center, font=\small}]
  \node[w, fill=blue!8] (s) at (0,0) {\zh{他}};
  \node[w, fill=gray!15] (v1) at (1.3,0) {\zh{说}};
  \node[w, fill=yellow!25] (o) at (2.7,0) {\zh{中文}};
  \node[w, fill=orange!20] (v2) at (4.1,0) {\zh{说}};
  \node[w, fill=red!15] (de) at (5.2,0) {\zh{得}};
  \node[w, fill=green!15] (adj) at (6.6,0) {\zh{很流利}};
  \node[font=\scriptsize, below=0.1cm of v1] {verb};
  \node[font=\scriptsize, below=0.1cm of o] {object};
  \node[font=\scriptsize, below=0.1cm of v2] {verb again};
  \node[font=\scriptsize, below=0.1cm of adj] {complement};
  \node[font=\small, anchor=west] at (7.7,0) {(a) repeat the verb};
  \node[w, fill=blue!8] at (0,-1.6) {\zh{他}};
  \node[w, fill=yellow!25] at (1.6,-1.6) {\zh{中文}};
  \node[w, fill=orange!20] at (3.0,-1.6) {\zh{说}};
  \node[w, fill=red!15] at (4.1,-1.6) {\zh{得}};
  \node[w, fill=green!15] at (5.5,-1.6) {\zh{很流利}};
  \node[font=\small, anchor=west] at (7.7,-1.6) {(b) object before the verb};
  \node[font=\small] at (3.3,-2.5) {Both: ``He speaks Chinese fluently.'' (\zh{流利} liúlì = fluent)};
\end{tikzpicture}
\end{center}

\begin{important}[Forms of the degree complement]
\begin{itemize}
  \item Negative: put \zh{不} before the \emph{adjective}, not the verb:
        \zh{我写得不好}, not *\zh{我不写得好}.
  \item Question: \zh{你写得好不好？} or \zh{你写得怎么样？} or \zh{你写得好吗？}
  \item Like adjective predicates, the plain adjective sounds contrastive, so \zh{很} is
        normally added: \zh{她唱得很好}.
  \item *\zh{他说中文得很好} is the classic error: \zh{得} cannot follow the object.
\end{itemize}
\end{important}

\begin{aside}
Three different particles are all pronounced \emph{de}, and Chinese students mix them
up too:
\zh{的} links a modifier to a noun (\zh{我的书}); \zh{得} links a verb to a complement
(\zh{说得好}); \zh{地} links an adverbial to a verb (\zh{慢慢地说}, ``speak slowly'').
In 101 we only need to produce \zh{的} and \zh{得}.
\end{aside}

\section{Other grammar}

\subsection{\zh{太\ldots{}了} and \zh{真}}

\zh{太} + adjective + \zh{了} means ``too'' or, in positive exclamations, ``so'':
\zh{太好了！} ``Great!'', \zh{太难了。} ``It's too hard.'' The negative \zh{不太} ``not
very'' takes no \zh{了}: \zh{语法不太难}. \zh{真} ``really'' expresses the speaker's
feeling: \zh{你的汉字写得真漂亮。}

\subsection{\zh{怎么} ``how''}

\zh{怎么} + verb asks about method.
\begin{itemize}
  \item \zh{这个字怎么写？} How do you write this character?
  \item \zh{``Waterloo''中文怎么说？} How do you say ``Waterloo'' in Chinese?
\end{itemize}

\subsection{\zh{有一点儿} vs \zh{一点儿}}

\begin{center}
\begin{tabular}{lp{5cm}l}
\toprule
Form & Position & Example \\
\midrule
\zh{有(一)点儿} + adj & before adjective; usually negative/undesired & \zh{第九课的语法有点儿难。} \\
adj + \zh{一点儿} & after adjective; comparison / request & \zh{请说慢一点儿。} \\
\zh{一点儿} + noun & before noun; ``a little'' & \zh{喝一点儿水} \\
\bottomrule
\end{tabular}
\end{center}

\begin{important}[Common error]
*\zh{这个字一点儿难} is wrong. Use \zh{这个字有一点儿难}.
\end{important}

\subsection{\zh{早}, \zh{晚} as greetings and complements}

\zh{早！} or \zh{早上好！} is a morning greeting. As complements:
\zh{你今天来得真早。} ``You're early today.''

\section{Talking about studying}

\begin{example}[After class]\leavevmode
\begin{itemize}
  \item \textbf{A:} \zh{你昨天的考试考得怎么样？}
  \item \textbf{B:} \zh{考得不太好。语法不难，可是汉字太多了。}
  \item \textbf{A:} \zh{你平常怎么学汉字？}
  \item \textbf{B:} \zh{我每天写，可是写得很慢。}
  \item \textbf{A:} \zh{我有一个好方法：我用卡片。你听录音吗？}
  \item \textbf{B:} \zh{不常听。我的发音有一点儿问题，我得多听录音。}
\end{itemize}
(\zh{方法} fāngfǎ = method; \zh{卡片} kǎpiàn = flash card.)
\end{example}

\begin{aside}
Study tips the instructor repeated: write each new character at least ten times
while saying its sound and meaning aloud; learn characters through their components
(\zh{谢} = \zh{讠} speech + \zh{身} body + \zh{寸} inch); preview (\zh{预习}) the
recording before class and review (\zh{复习}) after. Spaced-repetition apps like Anki or
Pleco's flashcards work well. Traditional forms: \zht{學習} for \zh{学习},
\zht{語法} for \zh{语法}.
\end{aside}

\section*{Review}
\begin{itemize}
  \item \zh{第} + number = ordinal.
  \item Degree complement: V + \zh{得} + (\zh{很}) + adj; with an object repeat the verb or
        front the object; negate the adjective.
  \item \zh{的/得/地}: three \emph{de}s with different jobs.
  \item \zh{太\ldots{}了}; \zh{不太} without \zh{了}.
  \item \zh{怎么} + verb = how.
  \item \zh{有一点儿} + adj (undesirable) vs adj + \zh{一点儿} (comparison/request).
\end{itemize}
